\section{From buffer to stream}
\label{sec:frombuffer}

In \kapref{ch:frames} we assumed that a complete frame was sitting in a buffer. It never is.

A serial link delivers one byte at a time, whenever it arrives, and a node waiting for the next
byte must not stop and wait --- it has other things to do. Reception therefore has to be
\textbf{incremental}: take one byte, do whatever can be done with it, and return.

That is exactly what a state machine is for.

\subsection{The states}

One state per field, plus one meaning done:

\begin{console}
WaitForSof1 -> WaitForSof2 -> WaitForLen  -> WaitForType -> WaitForDst
            -> WaitForSrc  -> WaitForSeq1 -> WaitForSeq2 -> WaitForPayload
            -> WaitForChk1 -> WaitForChk2 -> Ready
\end{console}

The parser's public interface is small:

\begin{cppcode}
class Parser final
{
public:
    /** Feed the parser one byte. Returns false on a parse error. */
    bool processByte(std::uint8_t byte) noexcept;

    /** True when a complete frame has been framed. */
    bool isFrameReady() const noexcept;

    /** Validate and copy out the framed data. */
    bool extractFrame(Frame& frame) const noexcept;

    /** Discard any partial frame and start over. */
    void reset() noexcept;
};
\end{cppcode}

\subsection{Two responsibilities, deliberately kept apart}

Note that \code{isFrameReady()} and \code{extractFrame()} are two different things, and that this
is deliberate:

\begin{itemize}
  \item \code{isFrameReady()} says that the \textbf{framing} is done --- the right number of
    bytes have arrived, in the right places.
  \item \code{extractFrame()} says that the frame is \textbf{valid} --- SOF is right, the length
    is plausible, the type exists, and the checksum adds up.
\end{itemize}

A frame can be fully framed and still be broken. That is the normal case when something has gone
wrong on the wire, and that is why the two questions have to be asked separately.

\begin{cppcode}
bool Parser::extractFrame(Frame& frame) const noexcept
{
    if (!isFrameReady()) { return false; }

    // Deserialize validates SOF, payload length, frame type and checksum, so a frame that was
    // framed correctly but arrived corrupted is rejected here.
    return frame.deserialize(myBuf, static_cast<std::size_t>(myParsedBytes));
}
\end{cppcode}

The validation therefore lives in \code{Frame::deserialize()} and not in the state machine. That
is no accident: the state machine knows where the fields \emph{are}, the frame knows what they
\emph{mean}, and mixing the two responsibilities gives you a parser that has to change every time
the protocol does.

\begin{aside}
\note \code{extractFrame()} is \code{const} and therefore cannot reset the parser. That is a
deliberate choice: the caller knows whether it wants to try again or throw the frame away, and
calls \code{reset()} itself. The same pattern returns in the driver: a method reports, the caller
decides.
\end{aside}
